% Card: side checks (rook and pawn vs rook)

\diagram[side-checks]{4K3/4P1k1/8/8/8/8/r7/5R2 b - - 0 1}

White's pawn is one step from queening, and White's king stands in front of it. Yet with Black
to move it is a draw, thanks to
\vocab{side-checks}{side checks}{The defending rook checks the enemy king from the side, along
the ranks. The defence holds when the king cannot escape the checks without leaving its pawn}:

\moves{1...\R a8+ 2.\K d7 \R a7+ 3.\K d6 \R a6+ 4.\K c5 \R a5+}

and so on. White's king has nowhere to hide: as soon as it walks away from the pawn to escape
the checks, Black's king takes the pawn.

Side checks only work from a distance. There must be at least three files between the
defending rook and the pawn; closer than that, the attacking king walks towards the rook,
the checks run out, and the pawn queens.
